\documentclass[12pt]{article}
\usepackage{amsmath,amssymb,amsfonts}
\usepackage[margin=1in]{geometry}

\begin{document}

\section*{Chapter 7, Problem 2}

\textbf{Problem:} Show that
\[
\frac{d}{dx}\int_a^x f(s)\,ds = \int_a^x \frac{\partial f(s)}{\partial x}\,ds + f(x)
\]
when both sides of the equation are considered to operate on a function in an integrand; i.e., we really mean
\[
\int_{-\infty}^{\infty}\frac{d}{dx}\Bigl[\int_a^x f(s)\,ds\Bigr]\,g(x)\,dx = \int_{-\infty}^{\infty}\Bigl[\int_a^x \frac{\partial f(s)}{\partial x}\,ds + f(x)\Bigr]\,g(x)\,dx
\]
for any $f(x)$ with a continuous derivative.

\bigskip
\textbf{Solution:}

We use the Leibniz integral rule for differentiation under the integral sign with a variable limit. Consider the function:
\[
F(x) = \int_a^x f(s)\,ds.
\]

Here $f(s)$ may depend on $x$ as well (in the distributional/generalized-function sense relevant to Green's functions), so we write $f = f(s,x)$ for clarity. Then by Leibniz's rule:
\[
\frac{d}{dx}\int_a^x f(s,x)\,ds = f(x,x) + \int_a^x \frac{\partial f(s,x)}{\partial x}\,ds.
\]

This is the standard Leibniz rule, which holds when $f$ and $\partial f/\partial x$ are continuous. The first term $f(x,x)$ comes from differentiating the upper limit of integration, and the second term comes from differentiating the integrand.

To prove this in the distributional sense (as the problem requests), consider a test function $g(x)$ and examine:
\[
I = \int_{-\infty}^{\infty} \frac{d}{dx}\Bigl[\int_a^x f(s,x)\,ds\Bigr]\,g(x)\,dx.
\]

Integrate by parts (assuming $g$ has compact support or vanishes at infinity):
\[
I = -\int_{-\infty}^{\infty}\Bigl[\int_a^x f(s,x)\,ds\Bigr]\,g'(x)\,dx.
\]

Now we reverse the order of integration. The region of integration is $\{(s,x): a \le s \le x\}$, which we can write as $\{(s,x): s \ge a,\; x \ge s\}$. Therefore:
\[
I = -\int_a^{\infty}\int_s^{\infty} f(s,x)\,g'(x)\,dx\,ds.
\]

For the inner integral, integrate by parts in $x$:
\[
\int_s^{\infty} f(s,x)\,g'(x)\,dx = \bigl[f(s,x)g(x)\bigr]_s^{\infty} - \int_s^{\infty}\frac{\partial f(s,x)}{\partial x}\,g(x)\,dx.
\]

The boundary term at infinity vanishes (since $g$ has compact support), so:
\[
\int_s^{\infty} f(s,x)\,g'(x)\,dx = -f(s,s)\,g(s) - \int_s^{\infty}\frac{\partial f}{\partial x}\,g(x)\,dx.
\]

Substituting back:
\begin{align}
I &= -\int_a^{\infty}\Bigl[-f(s,s)\,g(s) - \int_s^{\infty}\frac{\partial f}{\partial x}\,g(x)\,dx\Bigr]\,ds \\
&= \int_a^{\infty} f(s,s)\,g(s)\,ds + \int_a^{\infty}\int_s^{\infty}\frac{\partial f(s,x)}{\partial x}\,g(x)\,dx\,ds.
\end{align}

Reversing the order of integration in the double integral back:
\[
\int_a^{\infty}\int_s^{\infty}\frac{\partial f}{\partial x}\,g(x)\,dx\,ds = \int_a^{\infty}\int_a^x \frac{\partial f(s,x)}{\partial x}\,ds\,g(x)\,dx.
\]

Therefore:
\[
I = \int_{-\infty}^{\infty} f(x,x)\,g(x)\,dx + \int_{-\infty}^{\infty}\Bigl[\int_a^x \frac{\partial f(s,x)}{\partial x}\,ds\Bigr]\,g(x)\,dx,
\]
which gives us:
\[
I = \int_{-\infty}^{\infty}\Bigl[f(x) + \int_a^x \frac{\partial f(s)}{\partial x}\,ds\Bigr]\,g(x)\,dx.
\]

Since this holds for every test function $g(x)$, we conclude that in the distributional sense:
\[
\frac{d}{dx}\int_a^x f(s)\,ds = f(x) + \int_a^x \frac{\partial f(s)}{\partial x}\,ds. \qquad \blacksquare
\]

\end{document}
